\chapter{Parameter Definitions and Conventions}
\label{ch:params}

Launch monitors from different vendors disagree partly because they measure
different things and partly because they \emph{define} things differently.
TrackMan's definitions~\cite{trackman40params,trackmanclubdata} are the
industry reference and are adopted here. The critical subtlety is
\emph{when} each quantity is defined: club-delivery parameters at the time
of \textbf{maximum compression}, ball parameters \textbf{immediately after
separation}.

\section{Club-delivery parameters}
\label{sec:clubparams}

\begin{table}[htbp]
\centering\small
\caption{Club-delivery parameters (TrackMan conventions). ``GC'' = the
geometric center of the club head.}
\label{tab:clubparams}
\begin{tabular}{@{}p{3.2cm}p{7.2cm}p{4.2cm}@{}}
\toprule
\textbf{Parameter} & \textbf{Definition} & \textbf{Notes} \\
\midrule
Club speed & Linear speed of the GC just prior to first contact &
  Not the impact-point speed; the toe moves up to $\sim$7\mph{} faster than
  the heel \\
Attack angle & Vertical direction of GC motion at maximum compression &
  $+$ = hitting up. PGA Tour driver avg $\approx-1.3\degs$;
  LPGA $\approx+3\degs$ \\
Club path & Horizontal direction of GC motion at maximum compression &
  $+$ = in-to-out (right of target for RH) \\
Face angle & Horizontal direction the face normal points, at the
  center-point of ball contact, at maximum compression & $+$ = open \\
Face to path & Face angle $-$ club path & Sign controls curvature \\
Dynamic loft & Vertical angle of the face normal at the contact point at
  maximum compression & Differs from static loft via shaft lean/bend, face
  roll, impact height \\
Spin loft & 3D angle between the club-motion direction (path, attack
  angle) and the face-normal direction (face angle, dynamic loft) &
  $\approx$ dynamic loft $-$ attack angle only when face-to-path
  $\approx 0$ \\
Swing plane & Vertical angle of the plane traced by GC motion vs.\ the
  horizon & \\
Swing direction & Horizontal angle of that plane's base vs.\ the target
  line & \\
Low point & Distance from GC at maximum compression to the lowest point of
  the swing arc & $+$ = low point ahead of the ball (ball-first contact) \\
Impact height / offset & Vertical / horizontal strike location relative to
  face center & Drives gear effect (\cref{sec:gear}) \\
Dynamic lie & Shaft angle vs.\ horizontal at impact & \\
Closure rate & Angular rate at which the face is closing in 3D
  (\si{\degree\per\second}) & Foresight GCQuad-class only \\
\bottomrule
\end{tabular}
\end{table}

Two definitional choices deserve emphasis because they are common sources of
inter-device disagreement:

\begin{enumerate}
\item \textbf{Which point on the club is tracked.} TrackMan defines club
speed at the \emph{geometric center} of the head and reconstructs that point
from the radar ``3D silhouette'' of the head~\cite{trackmanclubspeed}. A
naive Doppler processor instead locks to the strongest or fastest return ---
often the toe of a driver --- and reports a speed several mph high. This
single choice explains much of the chronic club-speed disagreement between
brands.
\item \textbf{Where on the face orientation is evaluated.} Face angle and
dynamic loft are defined at the \emph{contact point}, not the face center.
On a curved driver face (bulge and roll, \cref{sec:gear}) the local normal
at a toe strike differs from the center normal by several degrees, so
systems that measure face pose but not impact location are systematically
biased on off-center hits.
\end{enumerate}

\section{Ball-launch parameters}

\begin{table}[htbp]
\centering\small
\caption{Ball-launch parameters (defined immediately after separation).}
\label{tab:ballparams}
\begin{tabular}{@{}p{3.2cm}p{10.8cm}@{}}
\toprule
\textbf{Parameter} & \textbf{Definition} \\
\midrule
Ball speed & Speed of the ball's center of gravity at separation \\
Smash factor & Ball speed $\div$ club speed \\
Launch angle & Vertical takeoff angle vs.\ the horizon \\
Launch direction & Horizontal takeoff angle vs.\ the target line \\
Spin rate & Rotation rate about the (single) spin axis, in rpm \\
Spin axis & Tilt of the rotation axis relative to the horizon;
  $-$ = tilted left $\Rightarrow$ draw for a right-hander \\
Carry / side / total & Trajectory descriptors; carry is measured to the
  point at launch elevation \\
Apex, landing angle & Peak height; descent angle at landing \\
\bottomrule
\end{tabular}
\end{table}

A golf ball in flight has exactly one rotation vector
$\vect{\omega}$. ``Backspin'' and ``sidespin'' are components of that
vector, not separate spins. Systems that report sidespin (SkyTrak-style)
and systems that report spin axis (TrackMan-style) are related by
\begin{equation}
\label{eq:spincomponents}
S_{\mathrm{side}} = S \sin\theta_{\mathrm{axis}}, \qquad
S_{\mathrm{back}} = S \cos\theta_{\mathrm{axis}}, \qquad
\theta_{\mathrm{axis}} = \operatorname{atan2}\!\left(
  S_{\mathrm{side}}, S_{\mathrm{back}}\right).
\end{equation}
A useful rule of thumb from TrackMan's data: $1\degs$ of spin-axis tilt
produces roughly $0.7\%$ of carry as side-curve (about
$0.7$\,yd per 100\,yd)~\cite{perfectgolfswing}.

\section{The directness hierarchy}
\label{sec:hierarchy}

For any launch monitor, each reported parameter falls somewhere on a
hierarchy from directly measured to purely modeled. Anticipating the
detailed treatment in \cref{ch:radar,ch:camera}, the hierarchy for the two
sensing families is summarized in \cref{tab:hierarchy}. This table is the
skeleton of this whole review.

\begin{table}[htbp]
\centering\small
\caption{Measured vs.\ derived, by architecture. ``M'' = measured directly,
``D'' = derived through a physical model, ``E'' = estimated/model-fit,
``--'' = not available.}
\label{tab:hierarchy}
\begin{tabular}{@{}lcccc@{}}
\toprule
\textbf{Parameter} & \textbf{Radar (outdoor)} & \textbf{Radar (indoor)} &
\textbf{Camera (photometric)} & \textbf{Hybrid} \\
\midrule
Ball speed          & M & M & M & M \\
Launch angles       & M & M & M & M \\
Spin rate           & M$^{a}$ & M/E$^{a}$ & M$^{b}$ & M \\
Spin axis           & D$^{c}$ & E & M$^{b}$ & M \\
Carry               & M & E$^{d}$ & E$^{d}$ & E$^{d}$ \\
Club speed          & M$^{e}$ & M$^{e}$ & M$^{f}$ & M \\
Club path / attack  & M & M & M$^{f}$ & M \\
Face angle          & D$^{g}$ & D$^{g}$ & M$^{f}$ & M \\
Dynamic loft        & D$^{g}$ & D$^{g}$ & M$^{f}$ & M \\
Impact location     & -- & -- & M$^{f}$ & M$^{h}$ \\
\bottomrule
\end{tabular}

\smallskip
\raggedright\footnotesize
$^{a}$ Doppler harmonic sidebands; requires surface asymmetry and
sufficient flight (\cref{sec:radarspin}).
$^{b}$ Dimple-pattern registration (\cref{sec:dimplespin}).
$^{c}$ Inverted from trajectory curvature via the Magnus constraint
(\cref{sec:spinaxis}).
$^{d}$ Trajectory model integration from measured launch
(\cref{ch:flight}).
$^{e}$ Reference-point dependent; silhouette reconstruction on TrackMan.
$^{f}$ Requires face fiducials (Foresight/Garmin R50) or overhead
geometry (Uneekor/ProTee).
$^{g}$ D-plane inversion from ball launch + path on radar-only units;
optically assisted on OERT-class hardware (\cref{sec:faceangle}).
$^{h}$ Markerless via camera on TrackMan~4/iO.
\end{table}
